\documentclass[11pt,oneside]{book}

\usepackage[T1]{fontenc}
\usepackage[utf8]{inputenc}
\usepackage{libertinus}
\usepackage{iftex}
\ifPDFTeX
  \usepackage[protrusion=true,expansion=true,final]{microtype}
\else
  \usepackage[protrusion=true,expansion=false,final]{microtype}
\fi
\setlength{\emergencystretch}{2em}
\usepackage{geometry}
\geometry{
  paperwidth=6in,
  paperheight=9in,
  margin=0.92in,
  top=0.84in,
  bottom=0.9in,
  headheight=14pt,
  headsep=16pt,
}
\usepackage{parskip}
\setlength{\parskip}{0.55em}
\setlength{\parindent}{0pt}
\linespread{1.045}
\frenchspacing
\raggedbottom
\clubpenalty=400
\widowpenalty=400

\usepackage{xcolor}
\definecolor{ink}{HTML}{1C1917}
\definecolor{stamp}{HTML}{7C2D12}
\definecolor{quiet}{HTML}{78716C}
\definecolor{rule}{HTML}{D6D3D1}

\usepackage{titlesec}
\usepackage{fancyhdr}
\usepackage{needspace}
\usepackage[hidelinks,bookmarks=true,bookmarksnumbered=false]{hyperref}

\hypersetup{
  pdftitle={Instruction for a Door},
  pdfauthor={Vladimir Sumarov},
  pdfsubject={A fictional self-help book. The instructions are not advice.},
}

\titleformat{\chapter}[block]
  {\sffamily\bfseries\fontsize{22}{26}\selectfont\color{ink}}
  {}
  {0pt}
  {}
\titlespacing*{\chapter}{0pt}{0.2em}{0.9em}

\titleformat{\section}[block]
  {\sffamily\bfseries\scriptsize\color{stamp}}
  {}
  {0pt}
  {}
  [\vspace{0.28em}{\color{rule}\hrule height 0.4pt width 1.55em}]
\titlespacing*{\section}{0pt}{1.65em}{0.62em}

\pagestyle{fancy}
\fancyhf{}
\fancyhead[L]{\sffamily\scriptsize\color{quiet} Instruction for a Door}
\fancyhead[R]{\sffamily\scriptsize\color{quiet}\nouppercase{\rightmark}}
\fancyfoot[C]{\sffamily\scriptsize\color{quiet}\thepage}
\renewcommand{\headrulewidth}{0pt}
\fancypagestyle{plain}{%
  \fancyhf{}%
  \fancyfoot[C]{\sffamily\scriptsize\color{quiet}\thepage}%
  \renewcommand{\headrulewidth}{0pt}%
}

\setcounter{secnumdepth}{-1}
\setcounter{tocdepth}{1}
\renewcommand{\contentsname}{Contents}

\begin{document}
\color{ink}

% HTML-SKIP-BEGIN
\begin{titlepage}
\thispagestyle{empty}
\vspace*{0.32\textheight}
{\sffamily\bfseries\fontsize{36}{40}\selectfont Instruction\par}
\vspace{0.15em}
{\sffamily\bfseries\fontsize{36}{40}\selectfont for a Door\par}
\vspace{1.05em}
{\large
A fictional self-help book.\par}
\vfill
{\sffamily Vladimir Sumarov\par}
\vspace{0.35em}
{\sffamily\footnotesize\color{quiet} 2026\par}
\vspace{1.4em}
\end{titlepage}

\clearpage
\thispagestyle{empty}
\vspace*{\fill}
{\sffamily\footnotesize\color{quiet} Instruction for a Door\par}
\vspace{0.75em}
{\footnotesize
Copyright \textcopyright\ 2026 Vladimir Sumarov.\par}
\vspace{0.65em}
{\footnotesize
This is a work of fiction. Names, characters, places, and incidents
are either products of the author's imagination or are used
fictitiously. Any resemblance to actual persons or events is
coincidental. It includes strong language and sex.\par}
\vspace{0.65em}
{\footnotesize
Licensed under the MIT License.\par}
\vspace*{\fill}
\clearpage

\tableofcontents
\clearpage
% HTML-SKIP-END

\chapter*{The Procedure}
\markboth{The Procedure}{The Procedure}
\addcontentsline{toc}{chapter}{The Procedure}

First, become twenty-three.

Do this before reading any further. There is no point filling your head with advanced material if you won't do the basics.

You are in your car, eating something out of paper. The engine is off. You have parked facing away from the building so that, for the duration of your break, the building does not exist. An excellent start. You have already discovered the power of visualization.

On your phone, a man your age is explaining how he started with nothing.

He has better lighting than you. This is partly equipment and partly the absence of a sandwich resting against his cock. Move the sandwich. We are trying to improve your life.

``I was twenty-three,'' he says. ``No idea what I was doing. Working out of my parents' garage.''

There it is. The starting point. You can pause him while you get ready.

\needspace{4\baselineskip}
\section*{What you will need}
\markright{What you will need}
\addcontentsline{toc}{section}{What you will need}

A garage. It should belong to your parents. If you rent one yourself, you have already misunderstood the exercise.

The garage must contain a socket, a workbench, and a father who comes to the door to ask what you're making. He may be skeptical. Skepticism is useful later, when you tell the story. It gives you someone to prove wrong without requiring him to stop paying for the electricity.

You will also need your parents to remain married, or at least sufficiently coordinated to keep the garage. Do not get distracted by which one was right. We are using the building.

If you have the wrong parents, ring them now. Explain what you require. Try to sound constructive. There is no need to make them feel bad about something they can no longer change; that is my job, and I have several chapters to get through.

The man on your phone turns his camera toward the old workbench. It is still there. He runs his hand along the edge and laughs at a burn mark.

``First one nearly caught fire.''

He sells a little box that tells people when their equipment has stopped doing what it is supposed to do. You had assumed he sold courses about selling courses. Unfortunately, there is a box. He opens one. The inside looks as if somebody knew where everything should go before they put it there.

He explains a fault in the first version, then the change he made. You follow him. For a moment you forget to dislike his face.

He puts the lid back on.

``Probably shouldn't show you that. That's the entire business.''

You give a small laugh through your nose.

Steady. We have not established whether he deserves it.

Behind him, a woman reaches across the bench for a mug. Only her arm comes into view. He lifts the loose cable out of her way without looking around.

``Thanks, Mum.''

She takes a second mug from beside the camera.

``I just thought, fuck it,'' he says. ``What's the worst that can happen?''

An important question. Write it down. Leave plenty of room.

\needspace{4\baselineskip}
\section*{Step one: take the risk}
\markright{Step one: take the risk}
\addcontentsline{toc}{section}{Step one: take the risk}

Return to the garage at twenty-three. You have no children. Your back works without a preliminary discussion. You can sleep badly and still be a person the next morning. None of this counts as capital because nobody has found a tasteful way to put it on a watch.

Place your savings on the bench. Spend them on parts.

When those parts fail, spend the next savings. These should be available because you are still in the garage. If you are doing this exercise in a rented room and the next savings have gone toward remaining in the rented room, move back to the correct childhood and repeat.

The man says he lost everything on his first attempt.

We should respect that. For a while he had to get a job he didn't like. His mother turns on a tap somewhere off camera.

Now you do it.

Lose everything.

Not the dramatic version where you clear your desk and emerge from a taxi carrying a plant. You haven't got a plant at work. Put your finger on the banking app. Work out which payment would bounce first.

You don't open it. You know.

Good. Preparation complete.

\needspace{4\baselineskip}
\section*{Step two: be the sort of person who does this}
\markright{Step two: be the sort of person who does this}
\addcontentsline{toc}{section}{Step two: be the sort of person who does this}

The first self-help book I wrote needed a whole chapter to explain this step. I have since become more efficient.

Be clever enough to make the box. Be interested enough to make it after it fails. Be able to talk about it to a stranger without making the stranger wish he had taken the stairs. These are habits. Acquire them retrospectively.

If you were the sort of boy who took things apart, you are well placed. If you were the sort who took things apart and got hit, please try to remember the lesson in a more entrepreneurial way.

The man describes spending evenings at the bench. He liked it there. Put that down as discipline.

He points to a little scar on his thumb. He remembers exactly which stupid thing he did to get it. As he talks, he picks up a screwdriver and finishes closing the box. His hands keep working after his sentence has finished.

You used to make things.

Don't look up. I haven't said they were good.

There is a photograph somewhere of a thing you made. You took several photographs of it, actually, until you found an angle that didn't show the bad bit. Then you sent the best one to somebody and waited for them to ask how you'd done it.

They said nice.

He holds his box toward the camera. You can see the tiny gap at one corner where the lid hasn't seated. He sees it too. Stops talking, backs out the screw, tries again.

This time it closes.

I would have edited that out. We are supposed to be establishing authority.

\needspace{4\baselineskip}
\section*{Step three: start before you did}
\markright{Step three: start before you did}
\addcontentsline{toc}{section}{Step three: start before you did}

You have been staring at the little line beneath the video. It says he is your age.

You check it again. Perhaps it says the age he was when this was filmed. You open the description. You find the date. You do a small, private piece of arithmetic whose result you immediately try another way.

No. Your age.

I could have selected an older man. I have a folder full of them. Bald men beside boats, explaining patience. I have to think about the quality of the work.

The man brings up a photograph of his first batch of boxes. They are stacked on the garage floor. He is crouching beside them in shorts, younger than he was a moment ago.

You remember what you were doing then. Not all of it. Just enough.

An entire year arrives as one room.

You know the ceiling from the bed. The fridge starts. You wait until it stops.

You drag the video forward.

Too far. Now he is in a bigger building, and there are people whose wages he pays. You drag it back a little. He is telling a story about a customer who wouldn't answer the phone.

Stay there. We can use that.

``You get a lot of rejection,'' he says.

At last. An area in which you can offer him some fucking mentorship.

But he means a buyer who was on holiday. The buyer eventually called back. There was an order. He delivered it. The story has a terrible ending.

Try not to drag the video again. You will wear a hole through the bit where he was poor.

\needspace{4\baselineskip}
\section*{Step four: remove the excuses}
\markright{Step four: remove the excuses}
\addcontentsline{toc}{section}{Step four: remove the excuses}

You need paper for this step. The sandwich wrapper will do if you have eaten enough of the sandwich. Turn it grease-side down. You are allowed a practical victory.

List the reasons you cannot follow the procedure.

Wrong age. No garage. Rent. Whatever is going on with your teeth. Work that leaves you enough energy to watch another man work. A person you can't leave. Nobody you would have to tell if you did.

Choose your own. I dislike doing both sides of an exercise.

Now cross them out.

There. Look how little room they take up.

You can do this with a thumb across the phone if you haven't got a pen. I am flexible about methods. I have always said that the important thing is results.

On screen, his phone rings.

``Sorry. Just a sec.''

He looks at it, answers, and turns partly away from the camera. He asks what the customer can see. Then he listens. You hear a faint, irritated voice, but no words.

He reaches for the box he has just closed.

``I think I know. Hang on.''

He is unscrewing it again.

There is a cut in the video. When it comes back, the box is open, and he is writing something on the back of an envelope. He reads it back to the person on the phone, waits, then crosses out one line.

``Okay. I'll send you the other one. Keep that there for now.''

Another pause. He rolls the screwdriver between his fingers.

``Yeah. No, that's on me.''

He ends the call. He looks at the camera as if he has forgotten it was there.

``Sorry. Where was I?''

He could ask me. I had him practically finished.

He checks the envelope once more.

``Right. Taking action.''

\needspace{4\baselineskip}
\section*{Step five: take action}
\markright{Step five: take action}
\addcontentsline{toc}{section}{Step five: take action}

You could write to him.

A clean, short message. You could ask about the box. You understood more of it than you expected. There is a thing at work that makes you wonder whether something like it would fit.

You open the comments instead.

Someone has written \emph{Step one: rich parents.}

An efficient man. He has completed the whole book without the unpleasant intermediate stage of reading it.

You press the little thumb beneath his comment.

The number changes. Nothing else does.

The man in the video is standing in the garage doorway now. You can see part of the house behind him, through the glass. There is a coat hooked over a chair. Someone moves past it carrying the two mugs.

He says he owes his parents a lot.

That is careless of him. Several excellent comments will now require amendments.

He looks back into the garage, then into the camera.

``If you're thinking about doing something, just start. I wish I'd started earlier.''

You lock the phone.

For a moment the glass shows your face, held at the angle from which you normally look at other people. You lift it a little. Better.

The building is still behind the car. You wipe your fingers on the clean edge of the wrapper and fold the dirty part inward. There is a bite left, which you eat standing up with the door open.

You unlock the phone to check the time.

The video is where you left it. His mouth is open. Your thumb rests above it, and then you tap the little bookmark so you can find him later.

Preparation complete.

\chapter*{Whose Hour}
\markboth{Whose Hour}{Whose Hour}
\addcontentsline{toc}{chapter}{Whose Hour}

There are twenty-four hours in everyone's day.

I have checked. The number does not change because you work nights, or because your day began at the wrong end of someone else's. Twenty-four. You can write it down if the size of the number has been troubling you.

You are in the car again. Different day, same view of nothing. The building is behind you, where you left it. The video from last time is still saved. Do not open it.

On your phone, a woman with excellent skin is standing in a kitchen that has already been cleaned by a person we will meet at 6:15.

``People always ask where I find the time,'' she says.

She does not answer them. She smiles, which is quicker, and holds up a glass of water as if the question had been about water.

\needspace{4\baselineskip}
\section*{The morning}
\markright{The morning}
\addcontentsline{toc}{section}{The morning}

\textbf{5:00. Rise.}

No alarm. The body simply knows.

If your body does not know, teach it. This takes three weeks, or six, or until you stop telling people you are teaching it. Leave the phone in another room. The other room should not contain a person who will be woken by it, or who is already awake because a child has discovered the concept of morning. If it does, choose a different other room. Optimization begins with architecture.

\textbf{5:15. Water. Silence. Intention.}

She drinks the water while looking out of a window. There is a tree in the window. It is doing nothing, which she respects. She names three things she is grateful for. She does not say them aloud, because the camera is close enough to her mouth that we would be able to check.

You have a tree. It is on the far side of the lot, and it has a bag caught in it. You are grateful that the bag is not your problem. This is not one of the three things. Keep going.

\textbf{5:30. Movement.}

Forty minutes. She uses the floor, because equipment is a story you tell yourself about tomorrow. Her form is correct. A small dog watches her from a basket and does not ask to be let out, which means the dog has either been trained or has already been out.

You could do this in the parking space. The car beside you has a person in it, also not going inside. He is eating and watching his own rectangle of light. Two optimized men, facing opposite directions. I am proud of both of you, in the way a man is proud of a shelf he has not yet put up.

Do not wave at him. He is in his movement block. You would be an interruption, and interruptions are how ordinary people stay ordinary.

\textbf{6:15. The house, continued.}

She is showered. Her hair is doing a thing that hair does in the kind of bathroom where the water arrives at the temperature you asked for. She is at the counter with a notebook, and the notebook already has writing in it from a previous life in which she was also awake at 5:00.

A man crosses behind her carrying a school bag and a lunch that has been cut into shapes. He says nothing. She does not turn. The shapes are animals, or they were animals before the walk to the car. He pauses at the door.

``You've got the thing at four,'' he says.

``I moved it.''

``To when?''

``I'll send it.''

He leaves. The door makes the soft sound of a door that has never had to be kicked. Somewhere in the house a child asks a question. She answers it without looking up, and the answer is the right length. Not a speech. Not a no. The sound of a person who has answered that question on other mornings and expects to answer it on this one.

Write that down. \emph{Presence.}

You do not have a child asking you a question at 6:15. If you do, she is not in this video, and you are not at a counter with a notebook. You are in a car, and the question is waiting in a different building, with a person who has already texted you once and is pretending that was the only time.

Don't answer it yet. We are still in her morning.

\textbf{7:00. Deep work.}

Ninety minutes. The phone is in a drawer. She says this while holding the phone, which is a kind of miracle, and I want you to respect miracles when they occur in the service of a system.

No meetings. No messages. She has told her team that this block is sacred, and her team has agreed, because her team is a group of people whose mornings are also made of blocks, or because her team would like to remain her team. Both are acceptable. Alignment does not require you to inspect it.

During these ninety minutes she produces work that other people will later call effortless. You will not see the work. You will see her nod at a screen and then close it, satisfied, the way a person closes a screen when the hard part has already happened in a room you were not shown.

\textbf{8:30. Nourishment.}

Eggs. Something green. She eats at the table, not over the sink, because how you eat is how you live, and how you live is how you earn, and how you earn is why you are watching this instead of eating.

The man from the doorway is not at the table. The child is not at the table. The plates that are not hers have already been moved to the side of the sink, where they are not in the frame.

You have an egg. It is in a sandwich, and the sandwich is in paper, and the paper is on your thigh, where it has begun to leave a map. Eat it. Nourishment is nourishment. I am not a snob about vessels.

\textbf{9:00. The world.}

She opens the drawer. The phone has thirty-one notifications and none of them are emergencies, because emergencies have been trained to arrive at other hours, in other people's hands. She answers four messages in complete sentences. She declines a call by typing the word \emph{moving}, which is true in the sense that she is moving her thumb.

Then she looks into the camera and tells you that anyone can do this.

\needspace{4\baselineskip}
\section*{Anyone}
\markright{Anyone}
\addcontentsline{toc}{section}{Anyone}

I love this part. I could watch it the way other men watch a fire.

The man in the next car starts his engine. He has finished his rectangle. He looks at the building, then at his hands, then at the building again, and does the thing with his face that people do when the next place is not a surprise and still has to be entered. His morning block is over. Nobody filmed it. There will be no notebook.

You could follow him in. You have four minutes, which is a block if you are willing to call it one.

Your phone lights up on your thigh, under the sandwich paper. You don't have to turn it over to know the shape of the notification. It is the second one. The person who sent it has learned not to use your name in the preview, because the preview is a public place, and some questions are better as a gray line that could be anything.

\emph{You up?}

No. That was a different year.

This one says she can do the later slot if you can do the pickup, and that she already asked you, and that she is not angry, which is how some people spell the other thing when they are at work and the child is not yet a problem anyone else can see.

You were going to do the pickup. You had moved a thing of your own to make the hour. Then the shift changed, and the hour moved with it, and the hour you moved is now a hole with your name still attached to it in a message you have not opened.

Open it.

I am not going to tell you what to write. You know what to write. You have written it before, with different clocks in it. The clocks are the only part that changes. The rest is a man explaining that a system he does not control has once again produced a result he would like treated as weather.

She will say okay.

She will say it quickly, because she is in a room where saying it slowly would be a conversation, and she is not in a car facing away from the building. She has a counter, or a desk, or a stretch of floor that other people cross while she is still typing. Your message will arrive between two of their questions. She will move the pickup into a gap that used to be a gap. The gap will close. Something else, which did not have a name before you wrote, will not happen, and she will not text you about that part, because that part does not fit in the preview, and because she has already used up the sentence that says she is not angry.

\needspace{4\baselineskip}
\section*{The overrun}
\markright{The overrun}
\addcontentsline{toc}{section}{The overrun}

You were the overrun. Now you are the person who creates one. Progress. I don't say that to everyone.

You could refuse the pickup. You could also refuse the shift. You could stand in the lot and explain to the building that your morning was supposed to contain a block of silence and a tree. The building has heard this explanation. It keeps a list of people who have delivered it, and the list is not a schedule. It is a drawer.

Don't give the speech. Speeches are what men do when the hour has already been spent and they would like a receipt.

Type the okay.

Not her okay. Yours. Short. No clock in it if you can help it.

You type it. You delete the second sentence. You send the first sentence alone.

The little mark changes from sent to read.

She does not reply. The pickup has a person now. The person is not you.

Anyone.

You look at the saved video. You do not open it. A man who makes a box can still lose an hour, and you do not need him to tell you that while you are trying to feel like a system.

The other car is gone. The bag is still in the tree. Your break has three minutes left, which is not a block, and you spend them sitting with the engine off, because the car is warmer than the room you are about to enter.

Preparation complete.

The building is still there. You get out. The hour goes on without you.

\chapter*{The Explanation That Finishes}
\markboth{The Explanation That Finishes}{The Explanation That Finishes}
\addcontentsline{toc}{chapter}{The Explanation That Finishes}

Choose an enemy you can reach from bed.

You will be spending a great deal of time together. Convenience matters. A man who requires you to put on trousers before hating him will eventually lose out to a woman with a podcast.

For this exercise, however, we have something better than a podcast. We have an email.

You are in the little room at work where people are taken to be disappointed without distracting production. There is a box of tissues on the table. The cellophane is still across the opening. Whoever furnished the room believes in distress but has not prepared for it.

Your manager turns the monitor toward you.

The company has bought places on a training course. Paid time, a qualification, and a chance at work that doesn't involve waiting for someone else to decide you are ready. You asked for a place. You were told to put yourself forward. You did.

The email says the first intake is reserved for women and staff from minority ethnic backgrounds.

You read it again. It has not developed an asterisk.

``Right,'' you say.

Your manager nods as if you have agreed on a sensible route around some roadworks.

``There'll be further opportunities.''

``When?''

``We haven't got dates for those.''

You look at the monitor. There are dates for these.

\needspace{4\baselineskip}
\section*{Establish the facts}
\markright{Establish the facts}
\addcontentsline{toc}{section}{Establish the facts}

You can print this one. I like that in a grievance.

Your manager says the programme is designed to remove barriers.

You make a small noise.

``Something you'd like to add?''

There is. Quite a lot, actually. Some of it has grammar.

You ask whether you would qualify if you were a woman. He says he isn't going to get into hypotheticals. A moment ago he was quite prepared to discuss a future intake with no date, budget, or existence. But a woman version of you has strained the limits of responsible speculation.

You ask what was wrong with your application.

``Nothing. It's the eligibility.''

Excellent. You have failed perfectly.

He pulls the monitor back toward himself. The meeting has reached the point at which he would prefer you to be an attitude problem.

You remain sitting.

``There is still a role for you in this,'' he says.

For a moment you look up.

One of the people going on the course will be away on Fridays. You know her work. They would like you to cover it.

He explains the arrangement. He has dates for that too.

The unopened tissues sit between you.

\needspace{4\baselineskip}
\section*{Save your work}
\markright{Save your work}
\addcontentsline{toc}{section}{Save your work}

You have the email on your phone. It was sent to you. You don't have to remember his exact words or persuade yourself that his face meant something. You can open the message while cleaning your teeth, and there it is again, still saying it.

There is a satisfaction in possessing something that bad.

I advise making a copy. Put it somewhere separate from the application. The application contains your explanation of why you wanted the place, and that is a much less agreeable document.

You spent an evening on it. You changed ``I think I could'' to ``I am confident that I can.'' You put the missing \emph{e} into development before anyone saw it. At one point you stood up and walked to the kitchen because you had written a sentence that sounded like a person who might be selected.

Keep that out of the screenshot.

Crop the top. Leave the eligibility paragraph and the bit about barriers. Send those to the man who will understand.

His reply arrives before you've put the phone down.

\emph{Fuck off.}

Not you. Them.

You send the Friday part.

This time he calls.

\needspace{4\baselineskip}
\section*{Find a second opinion}
\markright{Find a second opinion}
\addcontentsline{toc}{section}{Find a second opinion}

You're home now, sitting on the edge of the bed because the chair has things on it. A sock has come off inside one of your work shoes. You dig it out while he reads your message back to you, getting angrier at the same words you got angrier at, which is an excellent use of a telephone.

``They actually put that in writing?''

``That's what I said.''

``You're good enough to do her fucking job, though.''

You put the sock down.

There it is. Better than the line you had in the room. Shorter, too. You should have rung him from under the table.

He asks who got the places. You tell him the one you know. Neither of you says she can't do the work. You have both seen her do it. That would spoil the argument, and the argument doesn't need it. He repeats the bit about Fridays.

You laugh. Properly, this time. Loud enough to hear yourself in the room.

He has a story from his place. A different programme, a different manager, the same cheerful explanation of why a door is not a door when you are the one standing outside it. He has brought dates. You don't need them, but he gives them anyway.

For once, nobody tells you to work on how you come across.

Let him finish. This is considerably better than the breathing exercise I used to put here.

``I'd tell them where to stick the Fridays.''

``Yeah.''

A brief pause. You both know where you will be on Friday.

He says it again, with a better place to stick them.

\needspace{4\baselineskip}
\section*{Make it useful}
\markright{Make it useful}
\addcontentsline{toc}{section}{Make it useful}

An isolated grievance will get you through an evening. For lasting results, you need something that can handle the rest of the week.

Try \emph{they don't want us there}.

Say it once. It fits the course. Say it again. It fits the meeting, including the tissues. It is particularly good with the word \emph{barriers}. You can say that word in a voice now. Previously it was only a word.

Your friend says it with you.

The phrase is still working when you leave the course behind.

The company newsletter comes up. The photographs. The people who are supposed to look like the future of the place. The two of you can name the men who keep it running, although none of those men appears to own the correct kind of shirt for a photograph.

Then the job advertisements. The ones asking for your story. You have a story. Apparently the part where you turn up every day is a synopsis.

He says the country doesn't give a shit about men like you.

You take off the other shoe.

You mention a woman who told you to check your privilege. It was online. You remember her avatar more clearly than anything she said after that. Your friend has met her too, or someone who uses the same phrases and owns the same opinion of you. There seem to be a lot of her. Neither of you has met the woman who is meant to want a man who has been informed of all this.

I could give you a whole chapter on privilege. I could have you apologising to the chair by the end of it.

Instead, you go back to the Fridays.

The second telling is smoother. You leave out the little lift in your chest when the manager said there was still a role for you. It took no time at all to remove. In this version you knew what was coming.

Your friend says he would have walked out.

You say you nearly did.

Now we have two men who nearly walked out. An imposing force. Sit up a little.

\needspace{4\baselineskip}
\section*{Complete the explanation}
\markright{Complete the explanation}
\addcontentsline{toc}{section}{Complete the explanation}

``They don't want us anywhere,'' he says.

There. That is much easier to carry.

You don't have to open a job advert to know. You don't have to send the message and wait. No more sitting with your thumb over a button, wondering whether the wording sounds desperate. You can tell each other what the answer would have been. You're getting quite fast at it.

You have another example. It is old, and you have not thought about it for a while. At the time you couldn't tell whether the man had been an arsehole or whether you'd done something wrong.

Your friend settles it before you reach the end.

You thank him by settling one of his.

There is something almost luxurious about this. Usually you have to lie awake alone to conduct these proceedings. Tonight you have staff.

Your free hand has made a fist on your knee.

You open it. The nails have left marks.

The sock is still beside you. You pick a little hard flake off the heel and let it fall on the floor.

Your friend is still talking. You know what comes next, and you want to hear him say it anyway.

``You could do everything right.''

``Yep.''

``Wouldn't make a fucking difference.''

``Yep.''

He breathes into the phone. Somewhere at his end a cupboard closes.

I prefer this to optimism. Optimism needs something to happen.

For a while neither of you says anything. The line remains open. You can hear him moving about, putting things down. You move the sock into the shoe and push both shoes under the bed with your foot.

``You still there?'' he asks.

``Yeah.''

``Listen. I'll send you something.''

\needspace{4\baselineskip}
\section*{Supporting material}
\markright{Supporting material}
\addcontentsline{toc}{section}{Supporting material}

A screenshot arrives while he is still on the call.

You open it, ready.

It is a vacancy at his place. Different department. There is an early start in it, and a part about experience that you read twice.

``It's not that programme I was telling you about,'' he says. ``It's an actual job. Bloke's leaving.''

You scroll. There is no salary on the image.

``What's it pay?''

``Don't know the exact band. I'll ask.''

Of course he doesn't know. And you've heard about this place. From him. Extensively, and quite recently.

``Thought they were all cunts.''

``They are. He is, anyway. The one you'd be under.''

``You're really selling it.''

``He leaves you alone if the work's done. Mostly.''

You read the experience line again. Your friend doesn't speak over it.

Then another image arrives. An email he sent that morning, before you told him about the training. Your name is in it. Not just your first name. The full thing, spelled correctly, in the middle of a sentence saying you know what you're doing.

``You could have asked,'' you say.

``Yeah. Sorry.''

Below the message is the reply. The manager has asked whether your friend would be willing to work with you.

The answer is underneath.

\emph{Put him on with me.}

\chapter*{What the Payment Bought}
\markboth{What the Payment Bought}{What the Payment Bought}
\addcontentsline{toc}{chapter}{What the Payment Bought}

Improve the version of yourself that fits in an email.

The other one can wait in the car. He has practice.

Your friend has sent the band. It is more than you make now, and less than the number you say when someone asks what you should be on. There is a figure in the middle which he thinks they would actually start you at. He has put a question mark after it, which is honest of him and useless to you.

Early start means four. Not the sort of four you can approach from the correct direction if the evening has gone well. Four as in the day has not agreed to begin, and you will be there anyway, with your name on a list.

You read it twice. Then you do the sum.

The difference, after the earlier petrol and the worse week, is not a new life. It is a number that would stop one particular payment from arriving before the money does. You know which payment. You do not write its name down. Some improvements should be allowed to remain arithmetic.

``You'd be with me for the first month,'' he says. ``After that, his problem.''

``And the course?''

``What course?''

Quite.

\needspace{4\baselineskip}
\section*{Revise the applicant}
\markright{Revise the applicant}
\addcontentsline{toc}{section}{Revise the applicant}

You open the application from the training place. It is still confident. It still says \emph{development} with the \emph{e} in it. It is also still attached to a sentence about the future of the company, which will need to come out before it embarrasses anybody.

You delete the future.

You delete the line where you explained that you were ready for more responsibility. They have seen what you do with responsibility. It looks like Fridays.

In its place you put the work. Times, weights, the thing that broke and the order in which you noticed. You take out \emph{passionate}. You take out \emph{team player}. You leave in the year you spent on the machine nobody else wanted, because that year is ugly and specific and your friend says the man likes ugly and specific.

You read it back.

It sounds like a man who has already given up on being chosen and is now trying to be useful. I have sold books to that man. He is punctual, and he never asks what the ending is for.

Your friend reads it on his break.

``Don't say you were excluded.''

``I was.''

``Yeah. Don't say it. He'll think you're going to be a thing.''

You look at the sentence. It is accurate. It is also the first thing a stranger would use to decide what kind of trouble you are.

You cut it.

The application gets shorter and cleaner. Your friend says it sounds like you. This is not true, but it sounds like a person he can put in front of the man, which is a different achievement and should not be confused with the first.

You send it at 11:40 p.m., sitting on the bed, with one shoe off. The other shoe is still under the bed from last time. You do not retrieve it to make the moment more complete.

\needspace{4\baselineskip}
\section*{Attend in the correct voice}
\markright{Attend in the correct voice}
\addcontentsline{toc}{section}{Attend in the correct voice}

The interview is at 6:15, which is late for them and early for you. You park facing the building. This is a mistake. The building is a low brick shed with a roller door and a light that has been on all night. It does not reward being looked at.

Your friend meets you in the yard. He is already dirty in a way that means he has done something, not in the way that means he has dressed for a part.

``He'll ask about the early start,'' he says. ``Say you've done it.''

``I haven't.''

``Say you can.''

He looks at your shoes, then at your face, and decides not to improve either.

The man is older than both of you and does not shake hands. He has a mug with a crack in the glaze and a printout of your application, marked in pencil. He has underlined the year with the machine. He has not underlined anything else.

``Four o'clock,'' he says.

``I can do four.''

``Most people can do it for a week.''

You could tell him about the payment. You could tell him about the course, and the Friday cover, and the dates they haven't got. You have brought all of that with you. It is sitting behind your teeth, dressed for the occasion.

``I've done nights,'' you say. ``This is earlier. That's the difference.''

He makes a mark.

Your friend is standing by the door, not quite in the interview and not quite out of it. He does not help you. This is the correct behaviour, and you hate him for it with a purity you will have to thank him for later.

The man asks what you did when the machine went down.

You tell him. You leave out the part where you waited for someone senior to arrive and be pleased.

You give him the half hour.

He nods, once, at the paper.

``We'll let you know.''

\needspace{4\baselineskip}
\section*{Receive the result}
\markright{Receive the result}
\addcontentsline{toc}{section}{Receive the result}

They let you know on a Thursday, by email, while you are covering the Friday person's work a day early because she has a course day that is not a Friday. The sentence is short. The start date is in it. The figure is the one your friend put a question mark next to, without the question mark.

You read it in the car, on your break, with the engine off. The building you are about to walk back into is behind you. You have parked properly this time.

Your friend has already texted.

\emph{Told you.}

You type \emph{cheers}. You delete it. You type \emph{that figure?} He sends back the same number. You put the phone face down on your leg and look at the steering wheel for longer than the email requires.

It is a job. It is more money. It is a man who has met you and still written your name on a schedule. None of these is a mood. You could show them to someone who does not like you, and they would remain true.

You do not open the saved video. You do not send the email to your parents. You do not calculate what the number would do to a room, or a year, or the version of your life in which a person is waiting somewhere with her coat on. Those sums have a way of eating the result before you have had it.

The course email is still in the other folder. You do not move it. There are still no dates underneath it. The new job does not contain a qualification, a trainer, or a sentence in which anyone has decided you were the sort of person the first intake was for. You will be at work before the people on that course have woken up.

You write back to accept.

Your thumb hesitates over the word \emph{grateful}. You remove it. You remove \emph{excited}. You leave the date, the figure, and your name, spelled the way your friend spelled it when he was not asking your permission.

Then you sit with the door open and eat the rest of the sandwich, because the money has not arrived, and the shift has not started, and you are still a man on a break who was hungry before he became employed.

The email says you start on the 19th.

You swallow. You check the 19th. It is a Monday, which is considerate of it, and it is also eleven days away, which is long enough to tell one person and short enough to postpone telling the other.

You tell your friend first.

He says, ``Don't be late.''

You tell the woman second, in the same thread as the pickup, because that is the thread you have. You do not call it news. You write the start time and the money and ask whether she wants you back on the pickup.

She reads it.

Three minutes later she sends a time, an address, and the words \emph{if you're sure}.

You are not sure. You are employed at a figure that survives contact with the petrol and the payment, and you are awake at an hour you have not survived yet, and a person has not told you not to come. You type \emph{yes} and do not add a clock, and you do not add the sentence that would ask her to be pleased.

The phone goes dark.

Behind you, the building waits. In eleven days you will stop walking into it. Tonight you still have to. The sandwich paper goes into the door pocket. You get out, lock the car, and go back to the Friday work, which is still yours until the 19th.

The new email is in your pocket, unread by anyone you would like it to impress.

You do the covering anyway. The work is the same work. Your hands know it. At 4:10 you will know it in a different shed, under a man who does not shake hands, beside a friend who has already done an hour by the time you arrive.

\chapter*{Exit, with a Desk}
\markboth{Exit, with a Desk}{Exit, with a Desk}
\addcontentsline{toc}{chapter}{Exit, with a Desk}

Find your people.

Apparently they are behind a tyre place. You drove past twice, expecting alienation to have a better entrance.

It is Saturday. You have accepted the new job and have not started it. Your friend sent you the address with \emph{bring clothes you can fuck up}, which was considerate. You have several collections.

The photograph on the workshop's page showed a long bench and men who hadn't arranged their faces. Above it: \emph{No bosses. No bullshit. Just build.} You liked the page. Then you checked whether liking it was something people could see.

Now your friend is holding a door open with his foot.

``Thought you'd bottled it.''

``Couldn't find it.''

``It's where the address is.''

Inside, a man is trying to shut a metal toolbox. He presses the lid down with both hands. One end closes. The other rises.

``It's doing that to annoy you,'' your friend says.

``I know.''

He presses the other end. The first one rises.

``Have you tried the other end?''

``Have you tried fucking off?''

Your friend hangs up his coat.

The room smells of oil and something being heated that was once coffee. There are hand tools on the wall, gaps where other tools belong, and a broom leaning across a cupboard so that opening the cupboard requires acknowledging the broom. Somebody has drawn a cock in the dust on the window. It is an unambitious cock, with no programme for the future.

The man at the bench looks at your hands.

``Give us a hold of that end.''

You put your coat beside your friend's and hold the box.

\needspace{4\baselineskip}
\section*{Begin with what you have}
\markright{Begin with what you have}
\addcontentsline{toc}{section}{Begin with what you have}

The box is his. He reversed over it with his truck.

``Didn't see it,'' he says.

``What, behind you?'' your friend asks.

``I've had this conversation.''

One corner has folded inward. The tray is on the bench, still full of sockets. Under it is an old sandwich bag with a few screws in it and a receipt so pale that it might now be from any shop. He moves these out of the way carefully. The bag gets the same care as the tools.

You hold the box while he works on the corner with a block of wood and a mallet. After each attempt he tries the lid. Your friend watches with the expression of a man whose truck has never encountered an object he put behind it.

``Could buy a box,'' he says.

``Could buy mirrors,'' you say.

The man stops with the mallet raised. He looks at you. Your friend makes a noise and turns away.

``Brought another one, have you?''

He puts the block back against the corner.

``Hold it there, another one.''

You do.

The bench is at a height where you don't have to bend. Your phone is in your coat. The man asks you to move your thumb.

Move it. We should preserve your ability to scroll.

After a while your friend goes to make tea. You and the other man keep working. The lid begins to close, but the catch won't reach. He turns the box around. You catch the tray before it slides off the bench.

``Cheers.''

He says it to the hand that caught it. Then he is looking at the catch again.

Your friend asks whether you take sugar. He knows whether you take sugar. You have known each other for years. Apparently there are places where he still asks questions instead of forwarding evidence.

``One.''

``He's got no sugar,'' the man says.

``Then why did you ask?''

``Hospitality.''

You drink it without sugar. It is too hot and rather good, though I would prefer you not to build an identity around it.

For a while nobody says anything worth recording. Your friend finds the socket the man wanted earlier. It was underneath the box. This allows a previous argument to resume without either of them explaining it to you. You follow enough to enjoy the wrong man being right.

The tea cools. You have a little black grease in the fold beside your thumbnail. The catch now reaches, but only if you push the lid sideways. You try it twice, then show the man where it catches.

He leans in beside you to look.

You have been here nearly an hour. You have not checked whether anyone has read anything.

\needspace{4\baselineskip}
\section*{Establish common ground}
\markright{Establish common ground}
\addcontentsline{toc}{section}{Establish common ground}

This is where somebody should say what has happened to the country.

Instead, your friend asks about the truck.

``Still for sale?''

``Everything's for sale.''

``What do you want for it?''

``More than you offered.''

``You advertised it as a workhorse. It needs a fucking vet.''

The man checks the catch again.

``What are you using it for?''

``Taking stuff places. It's a truck.''

``I meant what's wrong with yours.''

``Nothing. I like yours.''

The man looks at him over the box.

Your friend drinks his tea.

That is the full inquiry. I would have made him do a worksheet.

The man tells him what he wants for the truck. Your friend says no. They discuss a scratch which your friend considers structural and the owner considers evidence of having used a truck. Neither changes his price. The truck remains where it was.

You turn the box a little. This time the lid seats without being pushed sideways.

``Try it,'' you say.

He opens it. Closes it. Works the catch.

Then he puts the tray back in and it won't close.

Your friend laughs so hard he has to put his mug down.

A socket is standing upright in the tray. You lay it on its side. The lid shuts.

The man opens it once more, slowly, while looking directly at your friend. He closes it. Click.

``I'll put that in the advert,'' your friend says. ``Comes with a box he can operate.''

The man picks up the box by its handle. Gives it a small swing. Puts it down again.

You want to come back.

I haven't asked you to. You have arrived at the conclusion by watching a man discover storage. My profession needs licensing.

\needspace{4\baselineskip}
\section*{Retain your independence}
\markright{Retain your independence}
\addcontentsline{toc}{section}{Retain your independence}

``Can you just turn up?'' you ask.

The man wipes his hands.

``First time you're his guest. After that you join. I'll show you.''

He walks to the far end of the room, where a narrow table sits under a noticeboard. There is a printer on it, mostly covered by a jacket. He takes off the jacket. Beneath it are a card reader, a box of pens, and a little plastic stand with \emph{MEMBERSHIP} inside it.

He turns the stand to face you.

Splendid. The free men have a desk.

You look back at your friend, who is washing the mugs with the solemn concentration of a man who has brought you somewhere.

``Thought there were no bosses.''

``There aren't,'' the man says. ``I'm the secretary.''

He has the form under your hand before you have worked out the difference.

Forty a month. A joining fee. Two evenings of supervised induction before you can book a bench; separate permission for the powered tools. Ordinary access is Tuesday evening and Saturday afternoon. A key requires approval, later. The key is not part of the forty.

``How much later?''

``Depends.''

An old friend. You had almost got through an afternoon without him.

The man shows you the cleaning rota. Members take turns. There is also a monthly meeting, which he says you should come to if you want a say.

``What about?''

``Last one was storage.''

``What about storage?''

``It went on a bit.''

On the wall behind him is the same photograph you saw online. The slogan is above it. Underneath, someone has pinned a sheet headed \emph{BENCH BOOKING PROCEDURE}. You can see the corner of a second sheet behind that one. The first sheet has replaced it without quite defeating it.

``Who approves the key?''

``Committee.''

``Who's on that?''

``Me. Him.'' He nods toward your friend. ``Couple of others.''

Your friend is drying a mug. You look at him until he feels it.

``What?''

``You're on a committee.''

``Somebody has to be.''

Listen to him. We have found the country.

The secretary asks whether you want the rules by email. You ask whether the monthly payment continues if you can't get there. He says yes. You ask whether you can keep things at a bench between visits. He points at the storage sheet, then looks tired enough that you believe in the meeting.

Forty would be possible when the new wages arrive. Today it is forty from the old wages. There is the joining fee as well. You stand looking at the card reader while he finds a pen that works.

Your friend comes over and leans against the wall. You don't want him to offer. You don't want him to move away and let you decide either.

``How many members?'' you ask.

``Enough to lose my good extension lead.''

There is a smaller notice beside the rules. \emph{DO NOT SHARE THE DOOR CODE.} It has been underlined twice.

You point at it.

``What happened?''

``Somebody brought a mate. Mate brought a mate. Next thing there's a bloke in here none of us knows.''

``What was he doing?''

``Using the place.''

You look past him at the wall of tools.

``Yeah, you can't have that.''

Your friend takes his hands out of his pockets.

The secretary puts the working pen in front of you.

\needspace{4\baselineskip}
\section*{Choose freely}
\markright{Choose freely}
\addcontentsline{toc}{section}{Choose freely}

You join.

I would have preferred a scene. A refusal on principle, ideally, delivered loudly enough that the tools knew why you were leaving.

Instead you pay the fee and the first month. The little machine asks whether you want a receipt. You do. You fold it and put it in your wallet rather than leaving it on the table where somebody might mistake you for a man who doesn't need receipts.

The secretary books your first induction for Tuesday. You check your current shift before agreeing. The second one will have to wait until you know how the new job leaves you. He won't take \emph{should be fine} for a booking.

``Text me when you know.''

He gives you his number. You ring it while standing beside him so he has yours. His phone vibrates under the jacket on the table. He digs it out, asks your name, and saves it. You put your phone back in your coat.

You go back to the bench. There is still a block of wood on it, the mallet, the tray's old sandwich bag. He takes the bag and puts it inside the toolbox. You pick up the mallet.

``Where's this live?''

Your friend points to a gap on the wall.

You put it there, then brush the flakes and dirt off the bench into a dustpan. The secretary shuts the box and carries it toward the door.

Your friend has his coat on. You get yours. The phone is still in the pocket. You put your hand around it and leave it there.

``Tuesday, then,'' the secretary says.

``Tuesday.''

Outside, your friend waits while you pull the door shut. You try the handle to check it has caught.

It opens.

``Harder,'' he says.

You pull it again, harder. It catches. You try the handle once more.

Your friend is looking at you.

``What?''

``Nothing.''

You give the door another pull before you follow him to the cars.

\chapter*{The Size of the Want}
\markboth{The Size of the Want}{The Size of the Want}
\addcontentsline{toc}{chapter}{The Size of the Want}

Have her round.

You will want a different room for this. That was an earlier exercise, and you skipped it. We are using the one you have.

She has sent nothing new. The time and the address are still sitting in the thread, above your yes, which you have reread often enough to know it by the shape. You could leave it there. A sensible man would arrive at her address when he said he would, do the pickup, and go home to a room she has not had to stand in.

You send your own address.

You type a sentence under it about the place being a bit of a state. You delete the sentence. You send the address by itself, like a man who has not noticed the state.

She reads it.

\emph{Now?}

\emph{If you want.}

\emph{Till ten.}

You look at the chair. The chair has been carrying your life for some time, and it was not hired for company.

\needspace{4\baselineskip}
\section*{Clear a surface}
\markright{Clear a surface}
\addcontentsline{toc}{section}{Clear a surface}

You move the work clothes into the bag. The bag looks like a person crouched in the corner. You push it behind the chair, where a person would also not want to find it.

The plate goes under the bed, next to the shoe. You eat the crisp that was on the plate. Hospitality.

There is a glass on the floor from last night, or the night before that. You take it to the bathroom, tip it out, and bring it back because a clean glass on the floor is not the improvement you thought it was. You put it in the wardrobe. The wardrobe already has a mug in it. You leave them together. They can explain each other.

You open the window. The room has a smell of its own, sleep and radiator dust and the towel on the door, and you would like her to meet a different one. Cold air comes in and does not take the job. You close the window before she can see you did anything about it.

You change your shirt. In the new one you look as if you were told a woman might come round. You change back. The first shirt knows the room. The room is less likely to give you away if you are dressed in it.

You sit on the bed. You stand up. You move the bag an inch deeper behind the chair and then stop, because a man who has prepared is a man who can be caught at it.

Your phone is on the chair. You leave it there. Face up. You are not waiting. The screen goes dark, and you let it.

She says nothing else. Ten minutes. You use them to take the towel off the back of the door and hang it on the back of the door again, smoother. Sit down.

The phone lights. Not a call. Her.

\emph{Here.}

You do not look in the mirror. The mirror has already done its work on other nights, and you know the angle you prefer. She will not be standing at that angle. You let her in.

\needspace{4\baselineskip}
\section*{Do not make it strange}
\markright{Do not make it strange}
\addcontentsline{toc}{section}{Do not make it strange}

She has the cold on her coat, and a bag on her shoulder. When you take the coat, the smell comes up out of the collar, weather and skin and a sweetness from whatever she drank on the way, and you know it. You have stood next to it in a car. You hold the coat a moment longer than a hook requires. There is a hook with a belt on it. You put the coat over the belt. She keeps the bag.

``Did you find it alright?''

``I found it,'' she says.

She looks at the chair. She looks at the bed. She keeps her shoes on for a minute, which you understand, and then she takes them off because her feet are cold and the radiator is the only hospitality you have managed that she can actually use.

``I've got till ten,'' she says. ``Mum's got him. I can't be late.''

``You won't be.''

``You don't know that.''

She is smiling. The smile is tired and a bit filthy.

You kiss her. The first one lands wrong, half on the corner of her mouth, and she puts her hand on your jaw and moves you the inch. Her fingers are cold. Her mouth is not. You make a sound you would not have authorised.

``Hello,'' she says.

You laugh, once, into her cheek.

She puts the bag by the bed and sits. The mattress slumps toward the wall. You should have warned her about the wall. You do not get the chance. She pulls you down by the first shirt, the one you were clever about.

Her neck smells more like the coat than the coat did. You stay there. She lets you. Her hand is in your hair, not steering yet, just there. You have a mouth. Use it.

You get her jeans down with the indignity that job always involves, one leg stuck, her heel on your thigh. She laughs. Not at you. At the leg.

There is a shoe under the bed. Your knee finds it. You shove it deeper in with the plate, a little cupboard of things you did not want in the evening, and you get your mouth on her.

She is already wet. This fact does not care about the wardrobe. You are so grateful you could become unpleasant about it. Don't. She has one hand over her face and then she takes it away.

You look. You look too much. She says your name once, low, annoyed, pleased. You keep the pace she sets with her heel on your back. When she comes she goes quiet first. The sound is after. It is small.

``Come here,'' she says.

You do. You push your jeans down. You are hard. She gets the condom out of her bag, not out of your drawer. She tears the packet with her teeth and rolls it on you without a ceremony, and the competence is so plain you nearly thank her.

``Don't,'' she says, though you haven't spoken.

The bed hits the wall. Then again. You try to hold the headboard and there isn't one, so you hold her hip and the sheet and a handful of nothing. She is warm and close and the room has nowhere left to put your attention. You last less than you meant to. Enough that you know it, and she knows it.

You drop your face into her shoulder.

``Don't,'' she says again, softer. She is not angry. She kisses the shoulder because your mouth is not available. Her lips stay there.

You stay inside her until you can't. She does not audit the interval. She breathes. Her hand moves once down your spine.

You deal with the condom in the bathroom. You wash your hands. You do not look in the mirror. When you come back she has the sheet over her and one foot still out. You get under the sheet. The wall is cold against your arm. She fits herself along you anyway.

Your phone lights on the chair.

You pick it up. Dad. Not a message. The call itself.

She feels you see it.

``Get it,'' she says. ``Might be something.''

``It's nobody.''

The screen does not agree with you. His name stays there, white on the dark. She looks at it in your hand. You watch her read it. Your thumb kills the light. You put it back on the chair.

``Nobody,'' she says. Better.

``He'll ring back if it matters.''

``Alright.''

She looks at the time. The looking changes her body before she speaks. You feel the leaving start in her shoulder.

``I have to go.''

``It's not ten.''

``I have to be back before I'm late, not when I am.''

You almost say stay. The word gets as far as your teeth, and you swallow it. The swallow is loud in your own ear.

``I'll do the pickup,'' you say.

``You already said.''

``I know.''

You say it again anyway, lighter. ``I'll be there.''

She sits up and finds her jeans. The leg that stuck before sticks again. She swears at it, fond of nothing, and you watch her stomach and the mark the sheet left on her side. You do not help until she asks.

``Shoe,'' she says.

You find hers by the radiator. You hold the back of it while she gets her heel in.

At the hook she puts her own coat on. You had wanted another go at the collar. There isn't one. She kisses you at the door, properly, no correction needed, and the kiss has a limit in it that the bed did not. She is already on the stairs in her head.

``Text me when you're in,'' you say.

``You'll be asleep.''

``I won't.''

She looks at you. She picks up the bag. She puts her hand on your chest, once, over the first shirt.

``Don't wait up like a mental person,'' she says. ``I'll text.''

Then she goes. You hear her on the stairs, and the front door, and her steps across the road.

You shut your own door.

The room is the same room. The sheet is pulled sideways. Her sweetness is on your mouth and on the sheet and already losing to the radiator dust.

You sit on the bed. You do not text.

You lie down where she was, with your face in the sheet, and breathe in until the smell is a fact you can keep.

\chapter*{Still Wanting, Still Here}
\markboth{Still Wanting, Still Here}{Still Wanting, Still Here}
\addcontentsline{toc}{chapter}{Still Wanting, Still Here}

Keep yourself ready.

There may be a life at short notice.

You are back in the room on another evening. The plate is out from under the bed. You have washed it, and the glass, and the mug from the wardrobe. A small amnesty. None of them has been asked to name the others.

The new job has not started. You have checked the route at the hour you will need it, laid out the clothes, and put the start date in your calendar. The calendar is pleased to have something with an address.

You put a tick beside \emph{clothes}. They are on the chair. They were on the chair before you made the list, but you are entitled to record work completed under an earlier administration.

There is another list on the phone. You have to scroll to find it. Sleep. Weight. Applications. The applications belong to a period when you were going to send three a day. A very organised period. You remember making the headings.

You add today's weight in the space beside yesterday's. The number is the same. Excellent retention.

Then you open the sleep page. Last night has not registered correctly. You know you were awake because you remember checking whether you were awake for longer than the night before. You correct the page. It now agrees that you slept badly.

You put the phone down and take your dinner out of the bag.

\needspace{4\baselineskip}
\section*{Maintain the equipment}
\markright{Maintain the equipment}
\addcontentsline{toc}{section}{Maintain the equipment}

You bought a pie on the way home. It is still hot through the paper, and the bottom has gone soft where the gravy leaked. You put it on the clean plate and sit on the bed with the plate on your knees. There is a fork. You went back for the fork.

The first piece of pastry breaks off too large. You eat it before attempting anything more architectural. The filling is peppery, with meat you have to cut against the side of the pie. You are hungry enough to stop holding your jaw in the position you use around other people.

For a while you eat.

I could have you log this. There is probably a pie in the database that you can use to misrepresent the one in your lap.

You get some gravy on your thumb. You lick it off and go after the piece caught in the fold of paper. The pie is good right down to the collapsed bit at the bottom. You finish it. You press the remaining pastry into the gravy with the back of the fork and eat that as well.

Your phone lights beside you. A reminder to record your evening meal.

You move the empty plate to the floor before returning your father's missed call.

He answers with your name, as if you have been called to a desk.

``Everything alright?''

``Yeah. You rang.''

``Did I?''

``The other night.''

``Oh. Yes.''

He takes a breath. You can hear him trying to recover a reason worth the call now being returned.

``Can't remember.''

You wait. He gives a little cough.

``Something on the television, I think.''

``Right.''

``You eating?''

``Just had a pie.''

``What sort?''

``Steak.''

``Good.''

He is more immediately pleased about the pie than you expected. You look at the plate. There is nothing you can tell him about it that you haven't told him already.

``I got another job,'' you say.

\needspace{4\baselineskip}
\section*{Bring your records up to date}
\markright{Bring your records up to date}
\addcontentsline{toc}{section}{Bring your records up to date}

You give him the date and the start time. He makes you repeat the start time.

``Four?''

``Four.''

``In the morning?''

``No, Dad. The other four before breakfast.''

He laughs. You put your heel on the edge of the bed and pull at a loose thread in your sock.

``More money?''

``Yeah.''

``How much more?''

You tell him. He goes quiet in order to do something with the number. You have already done it. You wait for him to finish.

``That'll help.''

``Yeah.''

``Is it what you wanted?''

You hook the thread around your finger. It comes away from the sock farther than you meant it to.

``It's more money.''

``I meant the work.''

``Same sort of thing. Different place.''

``Right.''

``Got everything you need?'' he asks.

``Clothes. Route. Alarm. I've checked it all.''

``Preparation complete,'' he says.

You hold the thread. Neither of you adds anything.

``Have to get some sleep, then,'' he says at last.

``Doesn't start yet.''

``No. When it does.''

``Yeah.''

He asks you to let him know how it goes. You say you will. Then there is the part of a call where both of you say variations of goodbye and neither hangs up, in case the other one has something.

He hasn't. You haven't. You end the call with the phone still against your ear.

On the screen, the meal reminder is waiting. You dismiss it. The sleep page is underneath, open at the night you have corrected.

You look around for something to do with the thread.

\needspace{4\baselineskip}
\section*{Make a claim}
\markright{Make a claim}
\addcontentsline{toc}{section}{Make a claim}

There is an envelope by the bed. Your name is on the front, spelled correctly, above a demand you have already paid. Keep the envelope. It is useful to have proof that the address works in one direction.

Turn it over. We shall use the back.

I have a form for this. I was going to put it at the end of the book, but you seem to have brought the matter forward by continuing to be here.

Copy the headings. Leave room underneath.

\textbf{CLAIM FOR THE LIFE REASONABLY EXPECTED}

\textbf{What was promised?} Attach the exact wording. If nobody said it in one go, reconstruct it from the school prospectus, your father's face, and the part of the advert after the car has stopped.

\textbf{What has been received?} List everything. Do not omit the job because you are annoyed with the question.

\textbf{What remains outstanding?} Use the space below.

\textbf{Who agreed to provide it?} One name, please. \emph{Everybody} is not a correspondence address.

\textbf{Acceptable replacement:} State what you would take instead.

\textbf{Signature:} Sign as the person who was meant to be here by now. If unavailable, sign yourself.

You write the new job under \emph{received}. The pen skips in the middle of the word. You go over it until it is considerably more employed than the rest of the form.

Under \emph{outstanding}, you start with \emph{a place}.

You stop. You already have a place. The envelope found it without assistance.

\emph{A home}, then. You cross out \emph{a place}. Under that you write \emph{someone to come home to}, which looks like a thing written by someone who has been told to write something. You nearly cross it out for that reason alone.

Leave it. I can read through the crossing-out.

You write \emph{kids} underneath. The word looks too casual, as if you were ordering something for the table. You press the pen into the full stop until the paper dents.

There are other things. You cannot get them into this corner without turning the envelope, and you don't want to turn it yet.

You skip the name of whoever owes you. You leave \emph{acceptable replacement} blank.

Then you sign.

Good. Your complaint now has exactly the same legal force as the back of an envelope, but it is much easier to store.

\needspace{4\baselineskip}
\section*{Retain a copy}
\markright{Retain a copy}
\addcontentsline{toc}{section}{Retain a copy}

You look at the form on your knees. You have written across the flap, so \emph{come home to} breaks when you lift it.

Press it flat.

We should establish a response time. Fourteen days sounds competent. Six weeks suggests an organisation with several floors. We can offer either, provided you understand that nobody has agreed to begin counting.

The pen is still in your hand. You put it down, then pick it up and try writing an address at the top. You get as far as \emph{To}.

Your father's number is in the recent calls. You can picture him answering again, more quickly this time. You put the pen down.

The room has cooled. You get the plate from the floor and put the fork on it. The thread from your sock is still wound around one finger. You unwind it into the pie paper and screw the paper into a ball.

You take the plate away and wash it. There is pastry stuck to the fork. You work at it with your thumbnail until it comes loose, rinse the fork, and bring both back.

The evening has not become late enough to go to bed. You check, in case.

On the phone is a saved video about using the last hour of the day. The man in the thumbnail is leaning against a kitchen counter with both hands around a mug. You can tell from his eyebrows that he has stopped doing something you still do.

You tap it. He says most people waste this hour.

You put him beside the envelope and go to the window. The glass is colder than you expect. You open it a little. Somewhere outside, a car door closes. You wait for the engine, then shut the window when the air starts coming round your ankles.

The man is still talking. You sit down and pull the bag toward you. The other parcel is still in the bottom.

An apple turnover, its paper dark with butter.

You move the envelope to the chair. You set the clean plate on your knees and take the turnover out carefully, keeping the sugar over the plate.

The man says you must learn to distinguish hunger from boredom.

You move him farther away so you can put your elbow down.

\chapter*{Irrelevant to the Claim}
\markboth{Irrelevant to the Claim}{Irrelevant to the Claim}
\addcontentsline{toc}{chapter}{Irrelevant to the Claim}

Skip the man who knows.

He is still on your phone, bookmarked, with his mouth open. You kept him so the exercise could fail with the materials already in the car.

The engine is off. The old building is behind you. The new one has not begun. Your break had a few minutes in it for eating, and you used those on the way here, standing at the counter, which means the car is only a seat and a screen.

\needspace{4\baselineskip}
\section*{Consult a substitute}
\markright{Consult a substitute}
\addcontentsline{toc}{section}{Consult a substitute}

You open the video at his mouth and drag it back to the call.

He asks what the customer can see. He listens. A faint voice comes out of the phone in his hand, still without words you can use. He unscrews his own lid. He writes on an envelope, crosses a line out, and says he will send the other one.

You have watched this often enough that your mouth wants the last line. Leave it. We are not starting a tribute.

There is a machine in the building behind you that stops, and nobody knows until someone walks over and puts a hand on it. You were not going to be taught the knowing. The dates for that are still missing. You leave his mouth where it is and open a different window. A blank line. You type the symptom the way you would tell a man on a break: it stops, and a person has to walk over.

The reply is there before you have decided how much of the machine you are willing to describe.

The first sentence is the one a careful man says. Isolate it. You already knew that. A phone that can recite the start of a warning has not become a colleague.

The second sentence wants the cover off. It has a place for you to look that the man in the video never showed, because he was opening his box and not yours. A third sentence starts a sequence. You do not read it. You have stood beside the machine. The cover stays on. A paragraph has not seen it, and you prefer the hands you currently have.

I had a better scene prepared if you continued. You would have learned about trust. I would have had the paragraph. You have spoiled it rather than be wrong in the building, which is a small reason, and not the one I wanted.

You put the phone on the passenger seat. The answer stays there, finished, face up, unused. His mouth is under it, still waiting, which is no way to leave a man you have been using.

\needspace{4\baselineskip}
\section*{Purchase the notice}
\markright{Purchase the notice}
\addcontentsline{toc}{section}{Purchase the notice}

The break is not over. You pick the phone up.

You are looking for the other object. Not his company. The part you understood: a notice when something stops using power.

The first result is a jaw. A clamp in a photograph, shut around one wire, sold as if the jaw were the knowing. Under the picture it asks for a board you do not have, and for a few lines of something the pleased reviewer already knew how to write. You are not that reviewer. You back out of him.

I would have bought the jaw. A small burn photographs better than a refusal, and you have denied me both the burn and the scene.

Under it is a plug. It goes into a wall. A thing goes into the plug. If the thing stops using power, a phone is told. There is an account. One complaint says it sends the notice when a kettle finishes and calls that a stop. Another says it stays quiet when the load has actually gone. You read them both. You understand them both, which is not a qualification and will not be accepted as one.

This will not reach the machine at work. That machine does not plug into your room. Your room is not becoming a demonstration.

The price is less than the forty you have already paid the men with the desk. It is also money you have. A large number can be refused with dignity. This one can be paid on a break, which is ruder.

You pay.

Your thumb stays on the glass longer than the bank needs. Above this payment, the workshop is still in the list. The new wages are not in the list. You watch the old ones thin by the price of a plug that has not yet arrived to do one job badly.

The screen says dispatched. Not to the car. To the room. The date is after the nineteenth. You will be in the other shed by then, under a man who does not shake hands, and the plug will be in a room that has never contained the machine.

There is no heading for this on the back of the envelope.

\needspace{4\baselineskip}
\section*{Do not file it}
\markright{Do not file it}
\addcontentsline{toc}{section}{Do not file it}

You do not open the other folder. The course email is still in it, waiting for dates. A plug in a parcel does not become the place you were refused.

You go back to the video. He has already said he will send the other one. You open the place where a message to him would go. A clean, short one. You understood the part where he listened before he took the lid off his own box. You have bought a worse notice. You have not fitted it to anything. The window waits.

You close it.

The answering line is still under the other window, ready to be believed if you scroll. You delete the thread. The warning goes with the sequence. You keep the order.

In the mirror the building is still there. You are not carrying the answer in. The machine will stop when it stops. A plug in a different room will have opinions about a kettle, later, if it arrives, if it connects, if you give it the account.

Your break ended while the payment was thinking. You are late. The phone is on your leg, lit with a thing you own and do not have. You lock it. His mouth goes. You put the coat on and put the phone in the pocket.

Your hand is on the door handle. The door is still shut.

\chapter*{Still Unsuccessful}
\markboth{Still Unsuccessful}{Still Unsuccessful}
\addcontentsline{toc}{chapter}{Still Unsuccessful}

First, become twenty-three.

You have had the whole book. I don't think a reminder is unreasonable.

You are sitting on the bed with one boot off. The new job has started. Four in the morning turned out to exist on the premises as well as in the email. Your friend was there. He showed you where to put your coat and which tap produced water you could drink. Nobody asked whether this was what you had pictured.

The money is still the agreed money. You haven't been paid it yet. You have, however, supplied a day for which it is owed. The room contains the boot you have taken off and the one you are negotiating with. It does not contain an audience.

You get a finger behind the heel and pull. The boot comes away with the sock inside it.

We can conduct the assessment barefoot.

\needspace{4\baselineskip}
\section*{Repeat the procedure}
\markright{Repeat the procedure}
\addcontentsline{toc}{section}{Repeat the procedure}

The garage should still belong to your parents.

You look around. There has been no garage. The chair has your clothes on it again. There is room for your foot between the bed and the bag if you turn it sideways. You turn it sideways.

Your phone is in the coat. You take it out and check the order. The plug is still coming. There is a picture of it to help you through the interval in which you have paid for an object but cannot drop it.

You put the phone on the bed. The envelope is on the chair, partly under a sleeve. You pull it out. Your name and address are on one side. The life is on the other. There is still no name after \emph{To}.

You set it on the floor beside the boots.

A garage acquired now would be a garage acquired now. I specified the earlier one. Please don't start buying sheds and expecting me to move the tick.

You reach for the sock inside the boot. It has caught on the heel. You pull it free, turn it the right way out, and put it beside the other one.

The body should be twenty-three too. A photograph of it will not do. You have photographs. In one of them you were convinced you had already left it too late. You could explain a few things to that little cunt. He would be disappointed by who had turned up to explain them.

No photograph tonight. The phone is already warm in your hand.

You open the saved video. It is still at the call. You let it play until the man has put his screwdriver down. Then you open the description beneath it and find the address for enquiries.

A question costs nothing to type. Sending it is extra.

\needspace{4\baselineskip}
\section*{Supply your evidence}
\markright{Supply your evidence}
\addcontentsline{toc}{section}{Supply your evidence}

You put his address into a new message. The subject line asks what it is about. You type \emph{Question about the box}.

This will distinguish you from the customers who wish to discuss his shorts.

In the body you write that you saw the video. He may have suspected this from the address supplied in the video, but it is important to begin where you are confident.

You write that you used to work with a machine that could stop without anyone noticing. It is now \emph{used to}. You look at those two words. You have been away from the place for one working day and already sound like someone who might know better than the people still there.

You delete \emph{used to work with} and put \emph{have worked with}.

Much better. Now it is experience.

You type the name of the place. You delete it. You can describe the problem without giving a stranger the name of somebody who has not asked for his opinion.

The next sentence takes longer.

You do not know what his box can connect to. You do not know what the machine already records or who would let him near it. The thing you bought will not go on it. You have enough information to be irritating and not enough to order anything. This is a recognised position in purchasing, but you are not in purchasing.

You type:

\emph{What would you need to know to say whether yours would work?}

Then you put the phone down.

That was one question. You are sitting as though you have submitted a thesis with an erection in it.

You get up to take the boots off the envelope. One has fallen sideways across it. There is a mark on the bill side. You wipe it with your thumb, make it wider, and leave it.

When you sit down again the message is still there.

You could send only the question. He could answer it. He could send a price you cannot afford, or a list you don't understand, or nothing. You know how to wait for nothing. You have the furniture.

Instead you add:

\emph{Do you ever work with anyone who finds places these could be used?}

You read it without moving your mouth.

There. A much better reason to be embarrassed.

\needspace{4\baselineskip}
\section*{Mark the result}
\markright{Mark the result}
\addcontentsline{toc}{section}{Mark the result}

Have you returned to the correct childhood? No.

Have you started before you did? No.

Have you become the man in the video? You have become a man sending him an email. I will need a separate column for that, and I haven't ruled one.

The new job can stay. I am not disputing the job. You got there at four, and you will be expected again. Nobody has sent a second email saying they mistook you for someone with a garage. Keep the figure too. It will be useful when the bill arrives.

It does not alter the date you started.

You rub the place on your ankle where the boot has been. There is a pale dent above the red mark. You pull the other trouser leg up to see if it matches. It does.

You have been at work. You have come back here. There are evenings when you would like somebody to look up as you come in, and you would like that to be so ordinary that she could be annoyed about something else.

She has not moved into the chair because you got a better number in an email. Nobody has.

The room is quiet enough that you can hear the phone vibrate against the mattress. You pick it up too quickly. A delivery update. The parcel will come tomorrow.

You dismiss it and return to the message.

You take \emph{ever} out of the last question. Then you put it back. Without it, you sound as if you have an appointment. With it, you sound as if you would like one.

I can arrange the report either way.

You scroll to the beginning. The subject is still \emph{Question about the box}. You now have two questions. You add an \emph{s}. This is the sort of thing you can finish.

The message has your name underneath it. Your own address is above it. You read both as though there might be a way of sending it without quite being you.

There isn't.

You send it.

\needspace{4\baselineskip}
\section*{Final assessment}
\markright{Final assessment}
\addcontentsline{toc}{section}{Final assessment}

Nothing happens in the room.

The envelope is still beside your boot. The mark on it has dried. Your second sock is inside the first one now; you put it there while reading, without deciding to be tidy.

You open the sent message to check that it has gone. It has. You find a sentence you would change. You close it before you can write an apology for the sentence.

He has not answered. The man can take a few minutes. He has a business to run, which was one of the things you held against him.

You have not told him the course story. You have not said how good you might have been if somebody had asked sooner. The question is in his inbox without either qualification. He can still ignore it. You have made that considerably easier by making it short.

The assessment is complete.

Still unsuccessful.

I can put \emph{potential} underneath if you want something your father can repeat.

You look at the first question again. The machine was round a corner from where a person usually stood. You didn't put that in. You had been thinking about what to call yourself.

There is a sheet in the bag with the new shift details printed on it. You take it out and turn it over. The back is empty. You get the pen from beside the bed and draw where the person stands, then where the machine is. You put the corner between them. You look at it and move the person a little farther back.

It is a poor drawing. I had hoped to finish on something more professional.

You photograph it. Your hand covers the corner.

You take the photograph again.

\end{document}
